% Part: first-order-logic
% Chapter: introduction
% Section: semantic-notions

\documentclass[../../../include/open-logic-section]{subfiles}

\begin{document}

\olfileid{fol}{int}{sem}

\olsection{Semantische begrippen}

Wanneer we kijken of $\Sat{M}{!A}[s]$ geldt, laat de toekenning~$s$
voor het gemak alle !!{variable}s een waarde krijgen. We gebruiken
echter alleen de waarden van de !!{variable}s die in~$!A$ staan. Nog
preciezer: alleen de waarden van de \emph{vrije} variabelen in~$!A$
doen ertoe. Dat zullen we bewijzen. !!^a{sentence} heeft geen vrije
variabelen. Daarom hangt de vraag of een structuur haar vervult niet
af van~$s$. Als $!A$ !!a{sentence} is, spreken we af dat
$\Sat{M}{!A}$ betekent: ``$\Sat{M}{!A}[s]$ voor alle~$s$.'' Dit geldt
precies als $\Sat{M}{!A}[s]$ voor minstens één~$s$ geldt. Voor
!!{formula}s hebben we toekenningen aan variabelen dus nodig. Voor
!!{sentence}s maakt de gekozen toekenning geen verschil.

Nu ``$\Sat{M}{!A}$'' is vastgelegd, weten we precies wat ``het geval''
en ``waar in'' bij !!{sentence}s van de eerste-ordelogica betekenen.
Uit deze definitie van $\Sat{M}{!A}$ kunnen we drie belangrijke
begrippen afleiden: geldigheid, semantisch gevolg en vervulbaarheid. !!^a{sentence} is geldig,
$\Entails !A$, als iedere !!{structure} haar waar maakt. We schrijven
$\Gamma \Entails !A$ als iedere !!{structure} die alle !!{sentence}s
uit~$\Gamma$ waar maakt, ook~$!A$ waar maakt. Dan is de conclusie een
semantisch gevolg van die verzameling. Een verzameling !!{sentence}s is
vervulbaar als er één !!{structure} is die alle zinnen in die
verzameling tegelijk waar maakt.

!!{formula}s zijn met inductie gedefinieerd. Ook de regels die bepalen
wanneer een formule waar is, volgen haar opbouw stap voor stap. Daarom
kunnen we inductie gebruiken om eigenschappen van de semantiek te
bewijzen en om de begrippen hierboven met elkaar te verbinden. We gaan
enkele van die eigenschappen verzamelen en bewijzen. Ze zijn op zichzelf interessant,
maar we hebben ze vooral later nodig wanneer we semantiek met
!!{derivation}ssystemen vergelijken. Eerst moeten we ook nog enkele
syntactische begrippen en bewerkingen precies definiëren. Ook dat doen
we met inductie.

\end{document}
